\begin{frame}{Using Upper Ontologies}
    \framesubtitle{Limitations and pitfalls}

    \AddToShipoutPictureFG*{
    \AtPageUpperLeft{\put(-5,-15){\makebox[\paperwidth][r]{Requirement 3}}}
    }%

    \begin{itemize}
        \item Upper ontologies enforce a specific perspective on the world
        \begin{itemize}
            \item This creates barriers for underrepresented perspectives of what, for example, the energy system 'is' and what its existence (or lack thereof) implies for people or other things with intrinsic value, and representing those perspectives in data
            \item This could be a criticism of any applied formal language (and therefore all structured data)
        \end{itemize}
        \pause
        \item Upper ontologies can get very intangible, causing domain experts to be unable to validate a model
        \begin{itemize}
            \item Explain to an asset manager that their inspection report is actually also a "generically dependent continuant" (BFO)...
            %\item I consider myself reasonably knowledgeable about this subject but I'm not sure this is a correct example
        \end{itemize}
        \pause
        \item Upper ontologies are often less formal than they seem to suggest
        \begin{itemize}
            \item nen-2660 considers an Organization to be a Physical Object
            \item I'm still on the fence whether an 'upper ontology' is the way to go, or if 'upper business terms' are a more accurate representation
        \end{itemize}
    \end{itemize}
\end{frame}
